\section{Conclusion and Future Work}\label{sec:conclusion} 

This paper introduced a comprehensive and practical architecture for an e-commerce chatbot that integrates goal-oriented question recommendation. By combining a product knowledge graph, a multi-agent system powered by large language models, and a user simulation environment for evaluation, the proposed framework addresses the challenge of balancing user engagement with business-driven objectives.

Experimental results validate the critical roles of both personalization and structured guidance. Incorporating user history significantly increases the relevance and acceptance of suggested questions, while the knowledge graph provides the logical scaffolding needed to guide conversations coherently and drive users toward target products.

Future directions include automating knowledge graph construction and maintenance, enabling cross-category recommendation, and incorporating real-time personalization techniques to further enhance the system’s adaptability and effectiveness.


% \section{Conclusion}
% Bài báo này đã trình bày thành công một kiến trúc toàn diện và khả thi cho chatbot thương mại điện tử tích hợp hệ thống gợi ý câu hỏi định hướng mục tiêu. Bằng cách kết hợp một đồ thị tri thức sản phẩm, một hệ thống đa tác tử dựa trên LLM, và một môi trường mô phỏng để đánh giá, nghiên cứu đã giải quyết được bài toán cân bằng giữa việc phục vụ người dùng và việc thực hiện mục tiêu kinh doanh.

% Kết quả thực nghiệm đã khẳng định vai trò không thể thiếu của cả hai thành phần cốt lõi: (1) Việc tích hợp lịch sử người dùng giúp tăng đáng kể sự hấp dẫn và tỉ lệ chấp nhận của các gợi ý. (2) Đồ thị tri thức là xương sống cho phép hệ thống dẫn dắt cuộc hội thoại một cách logic và tự nhiên, và là yếu tố quyết định để đạt được các mục tiêu kinh doanh đã định trước.

% Hướng phát triển trong tương lai sẽ tập trung vào việc tự động hóa quá trình xây dựng và cập nhật đồ thị tri thức, mở rộng khả năng gợi ý liên danh mục (cross-category), và tích hợp các phương pháp cá nhân hóa thời gian thực (real-time personalization) để làm cho hệ thống ngày càng thông minh và hiệu quả hơn.